\documentclass[11pt,a4paper]{article}
\usepackage[T1]{fontenc}
\usepackage[utf8]{inputenc}
\usepackage{lmodern,amsmath,amssymb}
\usepackage[margin=25mm]{geometry}
\usepackage[hidelinks]{hyperref}
\hypersetup{pdftitle={On typical configurations the flow truncation is controlled at any coupling; the whole obs},pdfauthor={navstokgap research working notes}}
\newcommand{\tightlist}{\setlength{\itemsep}{0pt}\setlength{\parskip}{0pt}}
\setlength{\emergencystretch}{3em}
\setcounter{secnumdepth}{0}
\begin{document}
\hypertarget{on-typical-configurations-the-flow-truncation-is-controlled-at-any-coupling-the-whole-obstruction-is-the-large-field-tail}{%
\section{On typical configurations the flow truncation is controlled at
any coupling; the whole obstruction is the large-field
tail}\label{on-typical-configurations-the-flow-truncation-is-controlled-at-any-coupling-the-whole-obstruction-is-the-large-field-tail}}

The condition \(t\|G\|_\infty\lesssim1\) of
\href{flow-jacobian-truncation-error.md}{the Jacobian note} can be
evaluated rather than assumed. In the free theory the flowed magnetic
field has \[\big\langle|b_t|^2\big\rangle=\frac{g^2}{16\pi^2t^2},
\qquad\text{so}\qquad t\,\big\langle|b_t|^2\big\rangle^{1/2}=\frac{g}{4\pi},\]
a pure number depending on the coupling alone and on no scale. Since the
truncation error at range \(R=\kappa\sqrt{8t}\) is
\(\exp[2t\|G\|_\infty-2\kappa^2]\), a \textbf{typical} configuration is
handled at any coupling by choosing
\[\kappa\ \gtrsim\ \sqrt{g/(4\pi)},\] that is, a truncation range longer
than the flow radius by the square root of the coupling. The obstruction
is therefore not the typical field strength but the \textbf{tail}:
\(\|G\|_\infty\) is governed by the largest fluctuation anywhere in the
volume, and for a Gaussian field it exceeds \(\eta/t\) with probability
of order \(e^{-c\eta^2/g^2}\) per correlation volume, small at weak
coupling and only polynomially controlled by the volume. In the
Euclidean framework that probability is a statement about the measure
and can be absorbed, which is precisely the large-field estimate of a
constructive renormalization group. In the Hamiltonian framework there
is no measure to integrate over, only a state, and the same regions must
be excluded by an operator statement. This locates the remaining
difficulty in one place: \textbf{the large-field tail, and the absence
of a Hamiltonian substitute for the probabilistic estimate that controls
it.} Free-theory values are exact; their use for the interacting theory
at scale \(\ell\) is the standard perturbative estimate and is labelled
as such. Constants explicit; nothing promoted.

\hypertarget{the-flowed-field-strength-in-the-free-theory}{%
\subsection{1. The flowed field strength in the free
theory}\label{the-flowed-field-strength-in-the-free-theory}}

From \href{flowed-bound-free-field.md}{the free-field note} §3, the
flowed magnetic two-point function at coincident points is
\[\big\langle b_{t,i}(0)b_{t,i}(0)\big\rangle
=\int\!\frac{d^3k}{(2\pi)^3}\sum_i\tilde G_{ii}(k)
=\int\!\frac{d^3k}{(2\pi)^3}\,\frac{g^2k}{2}\,e^{-2tk^2}\cdot2
=\frac{g^2}{2\pi^2}\int_0^\infty\!k^3e^{-2tk^2}dk,\] and
\(\int_0^\infty k^3e^{-\alpha k^2}dk=1/(2\alpha^2)\) with \(\alpha=2t\)
gives
\[\big\langle|b_t|^2\big\rangle=\frac{g^2}{2\pi^2}\cdot\frac1{8t^2}=\frac{g^2}{16\pi^2t^2},
\qquad \big\langle|b_t|^2\big\rangle^{1/2}=\frac{g}{4\pi\,t}.\] The flow
time is the only scale, so the dimensionless combination is
\[t\,\big\langle|b_t|^2\big\rangle^{1/2}=\frac{g}{4\pi},\] independent
of \(t\) and of the lattice spacing. In the interacting theory at scale
\(\ell=\sqrt{8t}\) the same estimate with the running coupling
\(g(\ell)\) is the leading perturbative statement.

\hypertarget{typical-configurations-need-only-a-longer-range}{%
\subsection{2. Typical configurations need only a longer
range}\label{typical-configurations-need-only-a-longer-range}}

Corollary 3 of \href{flow-jacobian-truncation-error.md}{the Jacobian
note} gives truncation error \(\exp[2t\|G\|_\infty-2\kappa^2]\) at range
\(R=\kappa\sqrt{8t}\). Replacing \(\|G\|_\infty\) by the typical value
of Section 1,
\[\varepsilon_{\rm typ}\ \sim\ \exp\Big[\frac{g}{2\pi}-2\kappa^2\Big],\]
so \(\varepsilon_{\rm typ}\le e^{-1}\) as soon as
\[\kappa^2\ \ge\ \frac{g}{4\pi}+\frac12,
\qquad\text{i.e.}\qquad R\ \ge\ \sqrt{8t}\,\sqrt{\tfrac{g}{4\pi}+\tfrac12}.\]
At \(g=1\) this is \(R\approx0.8\sqrt{8t}\); at \(g=4\pi\) it is
\(R\approx1.2\sqrt{8t}\); the range grows only as \(\sqrt g\).
\textbf{On typical configurations the flow truncation is therefore
controlled at every coupling}, at the price of an effective Hamiltonian
whose range exceeds the scale by a factor \(\sqrt{g/4\pi+1/2}\).

This removes the reading of the small-field condition as a weak-coupling
restriction. The condition \(t\|G\|_\infty\lesssim1\) is restrictive
because of the supremum, not because of the typical size.

\hypertarget{the-tail-is-the-obstruction}{%
\subsection{3. The tail is the
obstruction}\label{the-tail-is-the-obstruction}}

\(\|G\|_\infty\) is the largest value of the field strength anywhere in
the volume. For a Gaussian field of variance
\(\sigma^2=g^2/(16\pi^2t^2)\) the probability that \(|b_t|\) exceeds
\(\eta/t\) at a given point is
\[\mathbb P\Big(|b_t|>\frac\eta t\Big)\ \sim\ \exp\Big[-\frac{\eta^2}{2\sigma^2t^2}\Big]
=\exp\Big[-\frac{8\pi^2\eta^2}{g^2}\Big],\] and the number of
independent points at resolution \(\sqrt{8t}\) in a volume \(V\) is
\(V/(8t)^{3/2}\), so
\[\mathbb P\Big(\|b_t\|_\infty>\frac\eta t\Big)\ \lesssim\ \frac{V}{(8t)^{3/2}}\,
\exp\Big[-\frac{8\pi^2\eta^2}{g^2}\Big].\] Two readings.

\emph{Euclidean.} This is a statement about the Wilson measure, and the
exponential beats the volume factor as long as
\(\eta^2\gtrsim(g^2/8\pi^2)\log(V/(8t)^{3/2})\). Regions where it fails
are the large-field regions, they occupy a fraction of the volume that
is exponentially small in \(1/g^2\), and a construction handles them by
a separate argument whose cost is paid against that small probability.
This is the structure of Balaban's programme (Commun. Math. Phys. 122
(1989) 355 and the accompanying series; metadata level) and the
derivation above says why the split is forced: the truncation error is
exponential in the local field strength, so the configurations must be
sorted by it.

\emph{Hamiltonian.} Here the object is a state, and \(\|G\|_\infty\)
enters an operator bound. There is no integration over configurations in
which a rare region can be given small weight; the exponential factor
\(e^{2t\|G\|_\infty}\) multiplies an operator norm and must be
controlled uniformly, or the argument must be localized so that a bad
region affects only nearby terms. Localization of that kind is what the
Lieb--Robinson bound of \href{lieb-robinson-kogut-susskind.md}{the LR
note} supplies in principle, and turning it into a substitute for the
probabilistic estimate is the question this note isolates.

\hypertarget{consequences-for-the-route}{%
\subsection{4. Consequences for the
route}\label{consequences-for-the-route}}

\textbf{The window between the two methods is bounded and explicit.} The
cluster criterion of \href{strong-coupling-uniform-gap.md}{the T2 note}
holds for \(g\ge g_0=(48N^2/[(N^2-1)\beta_*])^{1/4}\); the flow
truncation on typical configurations holds at every coupling with range
factor \(\sqrt{g/4\pi+1/2}\), and the tail estimate degrades as
\(e^{-8\pi^2\eta^2/g^2}\). So the two do not fail in complementary
regimes of the coupling, as an earlier reading suggested: the flow step
is limited by the \textbf{tail at any coupling}, and the cluster step by
the coupling itself. A proof must therefore control the large-field tail
at every scale where the cluster criterion has not yet taken over, which
by \href{blocking-step-obstruction.md}{the blocking note} §5 is a number
of steps logarithmic in the scale ratio, of order thirty for \(SU(2)\)
from a bare coupling \(g_{\rm UV}^2=1/2\).

\textbf{The Euclidean framework is the right one for those steps.} The
suppression that makes the large-field regions harmless is a property of
the measure, and it has no Hamiltonian counterpart. This is a reason,
derived rather than conventional, why the constructive programme is
carried out in the Euclidean formulation, and why the Hamiltonian route
of this programme has reached its natural limit at exactly this point.

\hypertarget{consequence-for-state}{%
\subsection{5. Consequence for STATE}\label{consequence-for-state}}

The small-field condition is not a weak-coupling restriction: typical
configurations are handled at any coupling by lengthening the truncation
range as \(\sqrt g\). The obstruction is the supremum over the volume,
and its control is probabilistic, hence Euclidean. The Hamiltonian route
therefore terminates here in a well-defined way, and the two questions
it leaves are: whether a Lieb--Robinson localization can replace the
large-field probability estimate, and, on the Euclidean side, whether
the flow-truncation formulation of the step simplifies the existing
constructive treatment. Either is a substantial project; the first is
the one native to this programme.
\end{document}
